% A.69: explicit obstruction for specified finite-level local data.
\begin{statement}{A.69}{Proved}{A nonconstant joint fibre character and an obstructed new coordinate}
Put
\[
 c_1=0,\quad c_2=1,\quad c_3=-\frac{9563}{59024},
 \qquad R=\frac1{17},\qquad c=-\frac{11}{68},
 \qquad a_i=R-c_i.
 \tag{A.69.1}
\]
Let $C$ be the smooth projective compactification of the Kummer
cover with affine equations
\[
                 T-c_i=a_iX_i^4\qquad(1\le i\le3).
 \tag{A.69.2}
\]
It has the rational point $P_0=(R,1,1,1)$ and a point
$P_2\in C(\Q_2)$ with
\[
 T(P_2)=\frac{16}{17},\qquad X_1(P_2)=2,
 \quad X_2(P_2)=\frac12,
 \quad X_3(P_2)^4=\frac{4341}{869}.
 \tag{A.69.3}
\]
Every field factor of the fibre of (A.69.2) at $T=c$
contains $\Q(\mu_8)$.  Its joint quadratic character
corresponding to $\Q(\sqrt2)$ gives the unramified class
\[
        \mathcal A=
          \left(\frac{T-c}{R-c},\,2\right)
                              \in\operatorname{Br}C,
 \tag{A.69.4}
\]
with evaluations
\[
 \mathcal A(P_0)=0,\qquad
            \operatorname{inv}_2\mathcal A(P_2)=\frac12.
 \tag{A.69.5}
\]
In particular this common fibre character does not have constant
evaluation on the old Kummer curve.

Define ordinary local coordinates by $t_2=16/17$ and
$t_v=R$ at every other finite or infinite place of $\Q$.
The three global classes $a_i\in\Q^\times/(\Q^\times)^4$
restrict to the classes of $t_v-c_i$ at every place.
There is no global class in $\Q^\times/(\Q^\times)^2$
whose restrictions are the classes of $t_v-c$ at every place.
Thus these specified, everywhere ordinary finite-level data
cannot acquire the new coordinate even modulo two.
\end{statement}
\begin{proof}
We give the actual change of coordinates from small coefficients.
In a coordinate $t$ take the old punctures
\[
 d_1=\frac4{15},\qquad d_2=-\frac{64}{15},
                 \qquad d_3=\frac{869}{868},
 \tag{A.69.6}
\]
the base point $t=0$, and the new puncture $t=1$.
The affine map
\[
            T=\frac{t-4/15}{-68/15}
                      =\frac{4-15t}{68}
 \tag{A.69.7}
\]
sends these three old punctures to the $c_i$ of (A.69.1),
the base point to $R$, the new puncture to $c$, and $t=-4$
to $16/17$.  It multiplies each difference $t-d_i$ by
$-15/68$, so the old equations
$t-d_i=(-d_i)X_i^4$ become exactly (A.69.2).
It also gives the identity
\[
                         \frac{T-c}{R-c}=1-t.
 \tag{A.69.8}
\]

At the new puncture the three fibre constants are
\[
 D_i=\frac{c-c_i}{a_i},\qquad
 (D_1,D_2,D_3)=
             \left(-\frac{11}{4},\frac{79}{64},\frac1{869}\right),
 \qquad D_1D_2D_3=-\frac1{256}.
 \tag{A.69.9}
\]
The last equality uses $869=11\cdot79$.
In every field factor of the fibre algebra let $Z_i$ be
the images of its three root coordinates.  Then
\[
                         (4Z_1Z_2Z_3)^4=-1.
 \tag{A.69.10}
\]
Thus $4Z_1Z_2Z_3$ is a primitive eighth root of unity,
and that field contains $\Q(\mu_8)$ and $\sqrt2$.
This is an actual joint character of A.57.
It is not in the span of the three individual squareclasses:
each $D_i$ has even valuation at two, whereas $2$ has odd
valuation there.

The geometric Kummer classes of the three linear functions
$T-c_i$ are independent modulo fourth powers: their valuations
at the distinct points $c_i$ give three independent vectors.
Consequently (A.69.2) is geometrically integral and has degree
$4^3$ over the $T$-line.  Its normalization has a unique smooth
projective model.  The symbol in (A.69.4) can have residues only
over $T=c$ and infinity.  At a point over $c$, its character
restricts to zero by (A.69.10).  At a point $Q$ over infinity,
the equations give
\[
              v_Q(T-c_i)=v_Q(T)=4v_Q(X_i).
 \tag{A.69.11}
\]
Its pole order is therefore divisible by four and annihilates
the quadratic character.  The residue formula for cyclic
algebras and purity for a smooth curve show that (A.69.4)
belongs to $\operatorname{Br}C$
\cite[\S1.4 and Corollary 3.7.3]{AgentCTSkoro2019}.

At $t=0$ the three root coordinates are one.  At $t=-4$
their fourth powers are
\[
             16,\qquad\frac1{16},\qquad\frac{4341}{869}.
 \tag{A.69.12}
\]
The last number is a two-adic unit congruent to one modulo
sixteen, since both numerator and denominator are congruent
to five modulo sixteen.  Every element of $1+16\Z_2$ has
a fourth root in $1+4\Z_2$: the logarithm identifies these
two groups with $16\Z_2$ and $4\Z_2$, and division of the
logarithm by four followed by exponentiation gives the root.
This constructs the ordinary point $P_2$ in (A.69.3).
All its root coordinates are nonzero, and neither displayed
$T$ value is an old or new puncture.

At $P_0$ the first entry of (A.69.4) equals one, so its
evaluation is zero over every field.  At $P_2$ it equals
five by (A.69.8).  The exact dyadic Hilbert-symbol formula
for an odd unit $u$ is
\[
                    (u,2)_2=(-1)^{(u^2-1)/8}.
 \tag{A.69.13}
\]
For $u=5$ the exponent is three, so the quaternion algebra
is nonsplit and its invariant is $1/2$.
Here the needed nonsplitting can also be checked without
the general Hilbert-symbol formula.  If $5=x^2-2y^2$
with $x,y\in\Q_2$, the two summands have valuations of
different parity; valuation zero forces $x$ to be a unit
and $y$ to be integral.  Their norm is then $1$ or $7$
modulo eight, never $5$.  The quaternion norm criterion
\cite[Proposition 1.1.7]{AgentCTSkoro2019} proves the same
nonsplitting directly.
These two actual evaluations prove nonconstancy; unramifiedness
alone is not being used to infer it.

Finally, (A.69.2) at $P_0$ and $P_2$ proves the asserted
restriction identities for the three old classes modulo four.
Suppose a global squareclass $b$ had the indicated new local
values.  Divide it by the global constant $R-c=15/68$.
The quaternion algebra
\[
                       \left(b/(R-c),2\right)
                              \in\operatorname{Br}\Q
 \tag{A.69.14}
\]
would have invariant zero at every place other than two,
because there its first entry is locally a square.  At two
its invariant would be $1/2$, by (A.69.13).
This contradicts the global reciprocity sum for the Brauer
group \cite[Theorem 12.1.8]{AgentCTSkoro2019}.
It excludes the claimed global squareclass.
A completed Kummer class would have such a global reduction
modulo two, so it cannot give this prescribed new coordinate
either.
\end{proof}

\noindent\emph{Scope.}
The old data in A.69 are specified modulo four, with fixed
ordinary points at every place.  The old curve itself has the
rational point $P_0$; choosing that point at every place gives
different local data.  The proposition proves failure of a
new-coordinate extension for the displayed local tuple.
It does not turn that tuple into an ordinary completed pair,
and it does not disprove the all-rational-puncture version of
(G).  The distinction is recorded by the actual projection
from completed Kummer classes to their classes modulo four.

The certificate \path{code/joint_character_obstruction.py}
checks every rational coordinate and fibre identity above,
constructs the chosen fourth root modulo $2^{64}$, and checks
the exact Hilbert-symbol sign.  The infinite local root and
the reciprocity contradiction are proved in the text.
